\section{Discussion}\label{sec:discussion}

The failure of ThreadGate traces to a fundamental incompatibility between teacher-forcing and auxiliary attention losses, not to the gating architecture itself. Under teacher-forcing, $\mathcal{L}_\text{GAR}$ optimizes attention patterns conditioned on ground-truth prefixes that never appear at inference time. The gate parameters, subject to NTP gradient pressure that grows from 7.5$\times$ to 14.7$\times$ the GAR gradient magnitude (\S\ref{subsec:gradient-conflict}), converge toward 1.0 and become passive. The resulting 2--7\% residual attenuation cannot compensate for the distributional mismatch between teacher-forced and autoregressive contexts, as the 30$\times$ GAR gap confirms (\S\ref{subsec:dist-shift}). Stronger gating (lower initial values, higher $\lambda_\text{GAR}$) would degrade NTP performance, while weaker gating is already ineffective. This dilemma is structural: any lightweight intervention trained jointly with NTP under teacher-forcing faces the same pressure toward passivity.

Four directions follow from the experimental evidence. First, \textbf{inference-time attention intervention} sidesteps the training generalization problem entirely by rescaling or pruning attention heads at decoding time, operating on the actual autoregressive distribution rather than a teacher-forced proxy. Second, a \textbf{decoupled goal-processing pathway}, analogous to adapter modules but dedicated to goal tokens, could maintain goal information through a channel that NTP gradients do not dominate. Third, \textbf{GAR-aware decoding} would incorporate goal attention into the sampling or reranking stage rather than the training loss, avoiding exposure bias by construction. Fourth, \textbf{exposure bias mitigation for auxiliary losses} remains a viable direction: scheduled sampling \citep{bengio2015scheduled} or Professor Forcing \citep{lamb2016professor} applied specifically to $\mathcal{L}_\text{GAR}$ could reduce the distributional mismatch that destroys GAR transfer. Our v8 scheduled sampling experiment was invalidated by an implementation error; the approach itself warrants further study.

Our findings converge with \citet{dongre2026attention} from an independent angle. Dongre et al.\ showed that goal information persists in the residual stream but attention fails to route it to the output. If one were to attempt goal re-injection (directly inserting goal embeddings into the residual stream at late turns), the injected signal would likely face the same distributional mismatch we document: patterns learned under teacher-forced prefixes may not transfer to autoregressive contexts. Should both independent intervention routes (attention gating and goal re-injection) prove insufficient, this would suggest that lightweight modifications to the forward pass face a fundamental obstacle under NTP training pressure.

We scope our claims to lightweight residual gating combined with GAR auxiliary losses. More heavyweight interventions, such as full attention head retraining, separate goal encoders, or reinforcement learning from attention feedback, fall outside this investigation. The PCGrad diagnostic exclusion (\S\ref{subsec:pcgrad}) rules out gradient direction conflict but does not address magnitude imbalance \citep{chen2018gradnorm} or Nash equilibrium violations \citep{navon2022nash}, which remain untested. All experiments use a single model family (Qwen-2.5-7B-Instruct) on one dataset (UltraChat); the cross-model validation on Llama-3.1-8B-Instruct (\S\ref{sec:cross-model}) will determine whether these patterns generalize.